\input{preamble.tex}

\title{Отчет по лабораторной работе № 9. \\Использование протокола STP. Агрегирование каналов}
\author{Сущенко Алина Николаевна \\ НПИбд-01-23}
\date{2026}

\begin{document}

\maketitle
\newpage

\tableofcontents

\newpage
\section{Цель работы}
Изучение возможностей протокола STP и его модификаций по обеспечению отказоустойчивости сети, агрегированию интерфейсов и перераспределению нагрузки между ними.

\newpage
\section{Выполнение работы}

\subsection{Создание резервного соединения}
\begin{enumerate}
\item Заменили соединение между коммутаторами \texttt{msk-ansusenko-sw-01 (Gig0/2)} и \texttt{msk-ansusenko-sw-4 (Gig0/1)} на соединение между \texttt{msk-ansusenko-sw-01 (Gig0/2)} и \texttt{msk-ansusenko-sw-3 (Gig0/2)} (Рис. \ref{01}).

\begin{center}
    \centering
    \includegraphics[width=\textwidth]{../images/01.png}
    \captionof{figure}{Обновлённая топология сети.}
    \label{01}
\end{center}

\item На коммутаторе \texttt{msk-ansusenko-sw-01} настроили порт \texttt{Gig0/2} (Рис. \ref{02}):

\begin{center}
    \centering
    \includegraphics[width=\textwidth]{../images/02.png}
    \captionof{figure}{Настройка порта Gig0/2 на msk-ansusenko-sw-01.}
    \label{02}
\end{center}

\item На коммутаторе \texttt{msk-ansusenko-sw-4} настроили порт \texttt{Gig0/1} (Рис. \ref{03}):

\begin{center}
    \centering
    \includegraphics[width=\textwidth]{../images/03.png}
    \captionof{figure}{Настройка порта Gig0/1 на msk-ansusenko-sw-4.}
    \label{03}
\end{center}

\item На коммутаторе \texttt{msk-ansusenko-sw-3} настроили транковый режим на порту \texttt{Gig0/2} (Рис. \ref{04}):

\begin{center}
    \centering
    \includegraphics[width=\textwidth]{../images/04.png}
    \captionof{figure}{Ошибка при настройке транкового режима.}
    \label{04}
\end{center}

\item На коммутаторе \texttt{msk-ansusenko-sw-01} настроили транковый режим на порту \texttt{Fa0/23} (Рис. \ref{05}):

\begin{center}
    \centering
    \includegraphics[width=\textwidth]{../images/05.png}
    \captionof{figure}{Настройка транкового режима на порту Fa0/23.}
    \label{05}
\end{center}

\item То же самое сделали на \texttt{msk-ansusenko-sw-4} (Рис. \ref{06}):

\begin{center}
    \centering
    \includegraphics[width=\textwidth]{../images/06.png}
    \captionof{figure}{Настройка транков на msk-ansusenko-sw-4.}
    \label{06}
\end{center}

\item Проверили статус транков на \texttt{msk-ansusenko-sw-01} (Рис. \ref{07}):

\begin{center}
    \centering
    \includegraphics[width=\textwidth]{../images/07.png}
    \captionof{figure}{Проверка статуса транков на msk-ansusenko-sw-01.}
    \label{07}
\end{center}

\item Проверили статус транков на \texttt{msk-ansusenko-sw-3} (Рис. \ref{08}):

\begin{center}
    \centering
    \includegraphics[width=\textwidth]{../images/08.png}
    \captionof{figure}{Проверка статуса транков на msk-ansusenko-sw-3.}
    \label{08}
\end{center}

\item Проверили статус транков на \texttt{msk-ansusenko-sw-4} после настройки (Рис. \ref{09}):

\begin{center}
    \centering
    \includegraphics[width=\textwidth]{../images/09.png}
    \captionof{figure}{Проверка статуса транков на msk-ansusenko-sw-4.}
    \label{09}
\end{center}

\end{enumerate}

\subsection{Проверка связности с серверами}
\begin{enumerate}
\item С оконечного устройства \texttt{dk-ansusenko-01} проверили доступность file и web-серверов (Рис. \ref{10}):

\begin{center}
    \centering
    \includegraphics[width=\textwidth]{../images/10.png}
    \captionof{figure}{Пинг серверов с dk-ansusenko-01.}
    \label{10}
\end{center}

\item Убедились в правильном движении пактов ICMP в режиме симуляции через коммутатор \texttt{msk-ansusenko-sw-2} (Рис. \ref{11}):

\begin{center}
    \centering
    \includegraphics[width=\textwidth]{../images/11.png}
    \captionof{figure}{Пинг почтового сервера через коммутатор msk-ansusenko-sw-2.}
    \label{11}
\end{center}

\end{enumerate}

\subsection{Просмотр состояния протокола STP для VLAN 3}
\begin{enumerate}
\item На коммутаторе \texttt{msk-ansusenko-sw-2} просмотрели состояние протокола STP для VLAN 3 (Рис. \ref{12}):

\begin{center}
    \centering
    \includegraphics[width=\textwidth]{../images/12.png}
    \captionof{figure}{Состояние STP для VLAN 3 на msk-ansusenko-sw-2.}
    \label{12}
\end{center}

\item На коммутаторе \texttt{msk-ansusenko-sw-01} настроили его как корневой коммутатор STP (Рис. \ref{13}):

\begin{center}
    \centering
    \includegraphics[width=\textwidth]{../images/13.png}
    \captionof{figure}{Настройка корневого коммутатора.}
    \label{13}
\end{center}

\item Проверили состояние STP для VLAN 3 на \texttt{msk-ansusenko-sw-01} после настройки (Рис. \ref{14}):

\begin{center}
    \centering
    \includegraphics[width=\textwidth]{../images/14.png}
    \captionof{figure}{Состояние для VLAN 3 на корневом коммутаторе.}
    \label{14}
\end{center}

\item Убедились, что  пакеты ICMP пойдут от хоста dk-ansusenko-01 до mail через коммутаторы msk-ansusenko-sw-01 и msk-
ansusenko-sw-3. (Рис. \ref{15}):

\begin{center}
    \centering
    \includegraphics[width=\textwidth]{../images/15.png}
    \captionof{figure}{Движение пакетов через нужный коммутатор.}
    \label{15}
\end{center}

\item Убедились, что  пакеты ICMP пойдут от хоста dk-ansusenko-1 до web через коммутаторы msk-ansusenko-sw-01 и msk-ansusenko-sw-2. (Рис. \ref{16}):

\begin{center}
    \centering
    \includegraphics[width=\textwidth]{../images/16.png}
    \captionof{figure}{Верное движение пакетов.}
    \label{16}
\end{center}

\end{enumerate}

\subsection{Настройка Portfast}
\begin{enumerate}
\item На всех коммутаторах настроили режим Portfast на интерфейсах, к которым подключены серверы (Рис. \ref{17}, \ref{18}):

\begin{center}
    \centering
    \includegraphics[width=\textwidth]{../images/17.png}
    \captionof{figure}{Настройка Portfast на msk-ansusenko-sw-2.}
    \label{17}
\end{center}

\begin{center}
    \centering
    \includegraphics[width=\textwidth]{../images/18.png}
    \captionof{figure}{Настройка Portfast на msk-ansusenko-sw-3.}
    \label{18}
\end{center}

\end{enumerate}

\subsection{Изучение отказоустойчивости STP}
\begin{enumerate}
\item Запустили длительный пинг до почтового сервера с хоста \texttt{dk-ansusenko-01} с разрывом соединения. (Рис. \ref{19}):

\begin{center}
    \centering
    \includegraphics[width=\textwidth]{../images/19.png}
    \captionof{figure}{Длительный пинг до почтового сервера.}
    \label{19}
\end{center}

\item И без разрыва соединения. (Рис. \ref{20}):

\begin{center}
    \centering
    \includegraphics[width=\textwidth]{../images/20.png}
    \captionof{figure}{Длительный пинг до почтового сервера.}
    \label{20}
\end{center}

\item Сформируем агрегированное соединение интерфейсов Fa0/20 - Fa0/23 между коммутаторами msk-ansusenko-sw-1 и msk-ansusenko-sw-4 (Рис. \ref{21}):

\begin{center}
    \centering
    \includegraphics[width=\textwidth]{../images/21.png}
    \captionof{figure}{Логическая схема локальной сети с агрегированным соединением.}
    \label{21}
\end{center}

\end{enumerate}

\subsection{Агрегирование каналов EtherChannel}
\begin{enumerate}
\item На коммутаторах \texttt{msk-ansusenko-sw-1} и \texttt{msk-ansusenko-sw-4} настроили агрегированное соединение интерфейсов \texttt{Fa0/20 - Fa0/23} (Рис. \ref{22}, \ref{23}):

\begin{center}
    \centering
    \includegraphics[width=\textwidth]{../images/22.png}
    \captionof{figure}{Настройка EtherChannel для {msk-ansusenko-sw-1}.}
    \label{22}
\end{center}

\begin{center}
    \centering
    \includegraphics[width=\textwidth]{../images/23.png}
    \captionof{figure}{Настройка EtherChannel для {msk-ansusenko-sw-4} .}
    \label{23}
\end{center}

\end{enumerate}

\section{Выводы}
В ходе выполнения лабораторной работы были изучены протоколы STP и Rapid PVST+, а также технология агрегирования каналов EtherChannel. Настроено резервное соединение между коммутаторами, обеспечена отказоустойчивость сети, активирован режим Portfast для серверных портов. Проверена работоспособность ICMP-запросов в режиме симуляции. Все поставленные задачи выполнены. 

\end{itemize}

\end{document}
